\documentclass[11pt,a4paper]{article}
% preamble.tex — shared setup for human-verification documents
% Compile each document with:  tectonic <File>.tex

\usepackage[T1]{fontenc}
\usepackage[utf8]{inputenc}
\usepackage{lmodern}
\usepackage{amssymb,amsmath,amsthm}
\usepackage{mathtools}
\usepackage[dvipsnames]{xcolor}
\usepackage{listings}
\usepackage{hyperref}
\usepackage[margin=2.5cm]{geometry}
\usepackage{parskip}

% ── Colours ──────────────────────────────────────────────────────────
\definecolor{leanBg}{HTML}{F7F7F0}
\definecolor{haskellBg}{HTML}{F0F4F7}
\definecolor{keyword}{HTML}{7B3294}
\definecolor{comment}{HTML}{5A8C32}
\definecolor{noteFg}{HTML}{B35900}

% ── Lean 4 listings language ─────────────────────────────────────────
% NOTE: Lean listing blocks use ASCII approximations of unicode symbols
% (e.g. -> for arrow, o. for compose).  The literate key substitutes
% these with proper math-mode renderings in the PDF.
\lstdefinelanguage{Lean4}{
  morekeywords={def,theorem,structure,where,namespace,end,import,fun,
                let,have,show,by,exact,intro,cases,rfl,subst,obtain,
                funext,simp,rw,apply,variable,Type,Prop,instance,
                section,open,class,extends,deriving,sorry,noncomputable},
  sensitive=true,
  morecomment=[l]{--},
  morecomment=[n]{/-}{-/},
  morestring=[b]",
}

% ── Lean listing style ───────────────────────────────────────────────
\lstdefinestyle{lean}{
  language=Lean4,
  backgroundcolor=\color{leanBg},
  basicstyle=\small\ttfamily,
  keywordstyle=\color{keyword}\bfseries,
  commentstyle=\color{comment}\itshape,
  breaklines=true,
  frame=single,
  framerule=0.4pt,
  xleftmargin=4pt,
  xrightmargin=4pt,
  aboveskip=6pt,
  belowskip=4pt,
  literate=
    {->}{{$\to$}}2
    {=>}{{$\Rightarrow$}}2
    {<|}{{$\langle$}}1
    {|>}{{$\rangle$}}1
    {==>}{{$\Longrightarrow$}}3
    { o. }{{$\;\circ\;$}}3
    {forall}{{$\forall$}}1
    {:=}{{$\,:=\,$}}2
    {!=}{{$\neq$}}2
    {<=}{{$\leq$}}2
    {>=}{{$\geq$}}2,
}

% ── Haskell listing style ───────────────────────────────────────────
\lstdefinestyle{haskell}{
  language=Haskell,
  backgroundcolor=\color{haskellBg},
  basicstyle=\small\ttfamily,
  keywordstyle=\color{keyword}\bfseries,
  commentstyle=\color{comment}\itshape,
  breaklines=true,
  frame=single,
  framerule=0.4pt,
  xleftmargin=4pt,
  xrightmargin=4pt,
  aboveskip=6pt,
  belowskip=4pt,
  literate=
    {->}{{$\to$}}2
    {=>}{{$\Rightarrow$}}2,
}

% ── Theorem environments ────────────────────────────────────────────
\theoremstyle{definition}
\newtheorem{defn}{Definition}[section]
\newtheorem{thm}[defn]{Theorem}
\newtheorem{prop}[defn]{Proposition}

\theoremstyle{remark}
\newtheorem*{fidelitynote}{Fidelity note}

% ── Convenience macros ───────────────────────────────────────────────
\newcommand{\lean}[1]{\texttt{#1}}      % inline Lean identifier
\newcommand{\hask}[1]{\texttt{#1}}      % inline Haskell identifier
\newcommand{\Ext}{\mathrm{Ext}}
\newcommand{\Pos}{\mathrm{Pos}}
\newcommand{\Shape}{\mathrm{Shape}}


\title{Human Verification: \texttt{Containers.ProofSearch}\\[4pt]
  \large Binary Proof Search Trees as a Directed Container}
\author{Machine-generated from \texttt{lean/ProofSearch.lean} for review by Neil Ghani}
\date{July 2026}

\begin{document}
\maketitle

\begin{abstract}
This document presents every definition and theorem \emph{statement} from the
Lean~4 module \lean{Containers.ProofSearch} in three parallel forms: the
original Lean, a readable Haskell rendering, and precise mathematical English.
Proofs are omitted --- they are machine-checked by Lean.  The purpose is to
let a human verifier confirm that the formalisation captures the intended
mathematics.
\end{abstract}

\tableofcontents
\bigskip

%===================================================================
\section{Overview}
%===================================================================

This file realises binary proof-search trees as a directed container.  The
shapes are binary tree structures; positions are paths from the root to any
node.  The directed-container comonad structure comes from \emph{path
concatenation}:
\begin{itemize}
  \item $\mathrm{root}(t) = \varepsilon$ (the empty path \lean{here});
  \item $t \downarrow p = $ the sub-tree rooted at position $p$;
  \item $p \oplus q = $ the path obtained by following $p$ then $q$.
\end{itemize}

Laws D1 and D2 hold by $\mathrm{rfl}$ (definitionally from the first cases
of the recursive definitions).  D3--D5 are verified by structural recursion
on paths.

The induced comonad on $\mathrm{Ext}(\mathrm{proofSearch})$ has:
\begin{itemize}
  \item $\mathrm{extract}(t, v) = v(\varepsilon)$: read the label at the root;
  \item $\mathrm{duplicate}(t, v) = (t,\; p \mapsto (t \downarrow p,\; q \mapsto v(p \oplus q)))$:
        at every node, replace the label with the full sub-search from that point.
\end{itemize}
A coKleisli morphism $W(X) \to Y$ is a \textbf{proof strategy}: a function
that inspects the entire search tree to decide what to produce.

%===================================================================
\section{Shapes: binary search trees}
%===================================================================

% --- 1. STree -----------------------------------------------------
\subsection{Inductive type: \lean{STree}}

\begin{defn}[Binary search tree]
A \textbf{binary search tree shape} is either a \lean{leaf} (a terminal node)
or a \lean{node} with two sub-trees:
\[
  \mathrm{STree} \;::=\; \lean{leaf} \;\mid\; \lean{node}(\ell, r).
\]
A leaf represents either a dead end or a success state; the distinction lives
in the labelling, not the shape.
\end{defn}

\paragraph{Lean.}
\begin{lstlisting}[style=lean]
inductive STree where
  | leaf : STree
  | node : STree -> STree -> STree
\end{lstlisting}

\paragraph{Haskell.}
\begin{lstlisting}[style=haskell]
data STree = Leaf | Node STree STree
  deriving (Show, Eq)
\end{lstlisting}


%===================================================================
\section{Positions: paths in a tree}
%===================================================================

% --- 2. SPos ------------------------------------------------------
\subsection{Inductive type: \lean{SPos}}

\begin{defn}[Position in a binary search tree]
A \textbf{position} in an \lean{STree} is a path from the root to a node:
\begin{itemize}
  \item \lean{here}: the root (valid for any tree);
  \item \lean{left}$\,p$: descend into the left sub-tree and continue along $p$
        (only valid for a \lean{node});
  \item \lean{right}$\,p$: similarly for the right sub-tree.
\end{itemize}
For a \lean{leaf}, the only position is \lean{here} --- this is what makes a
leaf a dead end.
\end{defn}

\paragraph{Lean.}
\begin{lstlisting}[style=lean]
inductive SPos : STree -> Type where
  | here  : SPos t
  | left  : SPos l -> SPos (.node l r)
  | right : SPos r -> SPos (.node l r)
\end{lstlisting}

\paragraph{Haskell.}
\begin{lstlisting}[style=haskell]
data SPos t where
  Here  :: SPos t
  Left  :: SPos l -> SPos (Node l r)
  Right :: SPos r -> SPos (Node l r)
-- Requires GADTs or can be approximated as:
data SPos = Here | GoLeft SPos | GoRight SPos
\end{lstlisting}

\begin{fidelitynote}
In Lean, \lean{SPos} is indexed by the tree: \lean{SPos\,t} is the type of
positions \emph{in} $t$.  The \lean{left} constructor is only well-typed when
the tree is a \lean{node}.  In plain Haskell, the index is not enforced;
\lean{GoLeft Here} on a leaf would type-check but be semantically invalid.
The GADT version recovers the index.
\end{fidelitynote}


%===================================================================
\section{The directed-container operations}
%===================================================================

% --- 3. STree.sub -------------------------------------------------
\subsection{Definition: \lean{STree.sub}}

\begin{defn}[Sub-tree at a position]
The sub-tree at position $p$ in tree $t$:
\begin{align*}
  t \downarrow \lean{here} &= t, \\
  (\lean{node}\,\ell\,r) \downarrow \lean{left}(p) &= \ell \downarrow p, \\
  (\lean{node}\,\ell\,r) \downarrow \lean{right}(p) &= r \downarrow p.
\end{align*}
\end{defn}

\paragraph{Lean.}
\begin{lstlisting}[style=lean]
def STree.sub : (t : STree) -> SPos t -> STree
  | t, .here => t
  | .node l _, .left p => l.sub p
  | .node _ r, .right p => r.sub p
\end{lstlisting}

\paragraph{Haskell.}
\begin{lstlisting}[style=haskell]
sub :: STree -> SPos -> STree
sub t    Here        = t
sub (Node l _) (GoLeft  p) = sub l p
sub (Node _ r) (GoRight p) = sub r p
sub Leaf   _         = error "invalid position"
\end{lstlisting}


% --- 4. STree.shift -----------------------------------------------
\subsection{Definition: \lean{STree.shift}}

\begin{defn}[Shift (path concatenation)]
The shift $p \oplus q$ embeds position $q$ in the sub-tree at $p$ back
into the whole tree --- this is path concatenation:
\begin{align*}
  t.\lean{shift}(\lean{here}, q) &= q, \\
  (\lean{node}\,\ell\,r).\lean{shift}(\lean{left}(p), q)
    &= \lean{left}(\ell.\lean{shift}(p, q)), \\
  (\lean{node}\,\ell\,r).\lean{shift}(\lean{right}(p), q)
    &= \lean{right}(r.\lean{shift}(p, q)).
\end{align*}
\end{defn}

\paragraph{Lean.}
\begin{lstlisting}[style=lean]
def STree.shift : (t : STree) -> (p : SPos t) -> SPos (t.sub p) -> SPos t
  | _, .here, q => q
  | .node l _, .left p, q => .left (l.shift p q)
  | .node _ r, .right p, q => .right (r.shift p q)
\end{lstlisting}

\paragraph{Haskell.}
\begin{lstlisting}[style=haskell]
shift :: STree -> SPos -> SPos -> SPos
shift _          Here        q = q
shift (Node l _) (GoLeft  p) q = GoLeft  (shift l p q)
shift (Node _ r) (GoRight p) q = GoRight (shift r p q)
shift Leaf       _           _ = error "invalid position"
\end{lstlisting}


%===================================================================
\section{Directed-container laws}
%===================================================================

% --- 5. STree.shift_root (D3) ------------------------------------
\subsection{Theorem: \lean{STree.shift\_root} (D3)}

\begin{thm}[D3: right identity]
Shifting by the root of the sub-tree recovers the original position (right
unit law: $p \oplus \varepsilon = p$):
\[
  t.\lean{shift}(p, \lean{here}) \;=\; p
  \quad\text{for all } p : \mathrm{SPos}(t).
\]
\end{thm}

\paragraph{Lean.}
\begin{lstlisting}[style=lean]
def STree.shift_root : (t : STree) -> (p : SPos t) ->
    t.shift p .here = p
\end{lstlisting}

\paragraph{Haskell.}
\begin{lstlisting}[style=haskell]
-- D3 (right identity):
-- prop_shift_root t p = shift t p Here == p
\end{lstlisting}


% --- 6. STree.sub_shift (D4) -------------------------------------
\subsection{Theorem: \lean{STree.sub\_shift} (D4)}

\begin{thm}[D4: sub-trees nest along the shift]
The sub-tree at a concatenated path equals descending twice:
\[
  t \downarrow (p \oplus q) \;=\; (t \downarrow p) \downarrow q
  \quad\text{for all } p : \mathrm{SPos}(t),\; q : \mathrm{SPos}(t \downarrow p).
\]
\end{thm}

\paragraph{Lean.}
\begin{lstlisting}[style=lean]
def STree.sub_shift : (t : STree) -> (p : SPos t) -> (q : SPos (t.sub p)) ->
    t.sub (t.shift p q) = (t.sub p).sub q
\end{lstlisting}

\paragraph{Haskell.}
\begin{lstlisting}[style=haskell]
-- D4:
-- prop_sub_shift t p q = sub t (shift t p q) == sub (sub t p) q
\end{lstlisting}


% --- 7. STree.shift_assoc (D5) -----------------------------------
\subsection{Theorem: \lean{STree.shift\_assoc} (D5)}

\begin{thm}[D5: shift is associative]
Path concatenation is associative (with transport along D4):
\[
  t.\lean{shift}(p \oplus q,\; D4^{-1} \triangleright r)
  \;=\;
  t.\lean{shift}(p,\; (t \downarrow p).\lean{shift}(q, r))
\]
for all $p : \mathrm{SPos}(t)$, $q : \mathrm{SPos}(t \downarrow p)$,
$r : \mathrm{SPos}((t \downarrow p) \downarrow q)$.

The transport $D4^{-1} \triangleright r$ is needed because $r$ lives over
$(t \downarrow p) \downarrow q$, while the left-hand side needs it over
$t \downarrow (p \oplus q)$; D4 witnesses that these two types are equal.
\end{thm}

\paragraph{Lean.}
\begin{lstlisting}[style=lean]
def STree.shift_assoc : (t : STree) -> (p : SPos t) -> (q : SPos (t.sub p)) ->
    (r : SPos ((t.sub p).sub q)) ->
    t.shift (t.shift p q) ((t.sub_shift p q).symm [> r) =
      t.shift p ((t.sub p).shift q r)
\end{lstlisting}

\paragraph{Haskell.}
\begin{lstlisting}[style=haskell]
-- D5 (associativity, transport suppressed):
-- prop_shift_assoc t p q r =
--   shift t (shift t p q) r == shift t p (shift (sub t p) q r)
\end{lstlisting}

\begin{fidelitynote}
The Lean statement requires $r$ to be transported along $\mathrm{sub\_shift}$
(D4) from the fibre over $(t \downarrow p) \downarrow q$ to the fibre over
$t \downarrow (p \oplus q)$.  In Haskell, positions are unindexed and the
transport is invisible.
\end{fidelitynote}


%===================================================================
\section{The directed container}
%===================================================================

% --- 8. proofSearch -----------------------------------------------
\subsection{Definition: \lean{proofSearch}}

\begin{defn}[The \lean{proofSearch} directed container]
Define the directed container of binary proof-search trees as:
\[
  \mathrm{proofSearch} \;=\; (\mathrm{STree},\; \mathrm{SPos},\;
  \lean{here},\; \lean{STree.sub},\; \lean{STree.shift})
\]
with:
\begin{itemize}
  \item D1: $t \downarrow \lean{here} = t$ (by $\mathrm{rfl}$);
  \item D2: $t.\lean{shift}(\lean{here}, q) = q$ (by $\mathrm{rfl}$);
  \item D3: $\lean{STree.shift\_root}$;
  \item D4: $\lean{STree.sub\_shift}$;
  \item D5: $\lean{STree.shift\_assoc}$.
\end{itemize}
\end{defn}

\paragraph{Lean.}
\begin{lstlisting}[style=lean]
def proofSearch : DirectedContainer where
  Shape := STree
  Pos := SPos
  root := fun _ => .here
  sub := STree.sub
  shift := STree.shift
  d1 := fun _ => rfl
  d4 := STree.sub_shift
  d2 := fun _ _ => rfl
  d3 := STree.shift_root
  d5 := STree.shift_assoc
\end{lstlisting}

\paragraph{Haskell.}
\begin{lstlisting}[style=haskell]
proofSearch :: DirectedContainer STree SPos
proofSearch = DirectedContainer
  { shape = STree
  , pos   = SPos
  , root  = const Here
  , sub   = sub
  , shift = shift
  -- D1: sub t Here = t              (by definition)
  -- D2: shift t Here q = q          (by definition)
  -- D3: shift t p Here = p          (STree.shift_root)
  -- D4: sub t (shift t p q) = sub (sub t p) q  (STree.sub_shift)
  -- D5: shift t (shift t p q) r = shift t p (shift (sub t p) q r)  (D5)
  }
\end{lstlisting}


%===================================================================
\section{Translation conventions}
%===================================================================

\begin{enumerate}
\item \textbf{D1 and D2 are definitional.}  The first case in both
  \lean{STree.sub} and \lean{STree.shift} handles \lean{here}, giving
  $t \downarrow \lean{here} = t$ and $t.\lean{shift}(\lean{here}, q) = q$
  by $\mathrm{rfl}$.  D3--D5 require structural recursion.

\item \textbf{Positions are indexed.}  In Lean, \lean{SPos\,t} is the type of
  positions in $t$.  The tree structure constrains which constructors are
  valid.  Haskell approximates this with an unindexed \lean{SPos} type.

\item \textbf{The D5 transport.}  D5 requires a position to be transported
  along D4 before shift can consume it.  This transport is $\mathrm{rfl}$
  in both the \lean{here} and \lean{left}/\lean{right} cases (the recursion
  is definitional), but it must be stated in the type of D5 to make the law
  well-typed.

\item \textbf{Comonad reading.}  The induced comonad on
  $\mathrm{Ext}(\mathrm{proofSearch})$ is defined by the general machinery of
  directed containers; its laws follow from D1--D5.  This document only
  formalises the directed-container structure itself.

\item \textbf{Proof bodies omitted.}  Only statements are translated.
\end{enumerate}

\end{document}
